%% ma: L12-B18-0011
%% lop: 12
%% chuong: 6
%% bai: 18
%% dang: 12.18.7
%% loai: TN4
%% muc: TH
%% dap_an: "D"
%% nguon: "(NBV)-10 ĐỀ KIỂM TRA HK2-TOÁN 12-DUNG CHUNG 3 SÁCH.pdf (H01, Đề số 1), Phần 1 Câu 11"
%% kien_thuc: "Bài 18 (công thức xác suất toàn phần dạng P(A) = P(AB) + P(A B̄))"
%% trang_thai: cho-duyet
%% ly_do: "Dạng mới đề xuất 12.18.7 (khẳng định luôn đúng về phân tích một biến cố theo B và biến cố đối của B), chờ giáo viên duyệt."
%% ghi_chu: "Đáp án gốc D đúng. Đề gốc in lỗi điều kiện '0 < P(A) < 1, ) < P(B) < 1'; bản kho sửa thành 0 < P(A) < 1, 0 < P(B) < 1."
\begin{ex}
Cho các biến cố $A$ và $B$ sao cho $0<P(A)<1$, $0<P(B)<1$. Khẳng định nào sau đây luôn đúng?
\choice{$P(B)=P(AB)+P\left(A\overline B\right)$}{$P(A)=P\left(\overline A\,\overline B\right)+P(AB)$}{$P(B)=P\left(\overline AB\right)+P\left(A\overline B\right)$}{\True $P(A)=P(AB)+P\left(A\overline B\right)$}
\loigiai{Hai biến cố $AB$ và $A\overline B$ xung khắc và hợp của chúng là $A$ nên $P(A)=P(AB)+P\left(A\overline B\right)$ (D đúng). Tương tự $P(B)=P(AB)+P\left(\overline AB\right)$ nên A và C sai. Phương án B sai vì $AB$ và $\overline A\,\overline B$ không phủ hết $A$. Chọn D.}
\end{ex}
